\section{Sentencias DDL del Subsistema de Mensajería Privada (Sergio)}

A continuación se detallan las definiciones de las tablas utilizadas en el subsistema de Mensajería Privada. Este diseño garantiza la integridad de las comunicaciones bidireccionales y la persistencia de los estados de cada mensaje.

\subsection{Tabla: CONVERSA}
Esta tabla define la existencia de un canal habilitado entre dos usuarios. Su función es actuar como base de integridad referencial para la tabla de mensajes, validando que ambos participantes tienen permiso para interactuar entre sí.

\begin{lstlisting}[language=SQL, caption={Creación de la tabla CONVERSA}]
-- Creación de la tabla conversa
CREATE TABLE CONVERSA (
    IDUSUARIO1          NUMBER NOT NULL,
    IDUSUARIO2          NUMBER NOT NULL,
    CONSTRAINT PK_CONVERSA PRIMARY KEY (IDUSUARIO1, IDUSUARIO2),
    CONSTRAINT FK_CONV_U1 FOREIGN KEY (IDUSUARIO1) REFERENCES USUARIO(IDUSUARIO) ON DELETE CASCADE,
    CONSTRAINT FK_CONV_U2 FOREIGN KEY (IDUSUARIO2) REFERENCES USUARIO(IDUSUARIO) ON DELETE CASCADE
);
\end{lstlisting}

\textbf{Justificación de diseño:}
\begin{itemize}
    \item \textbf{Simetría en la Relación:} El diseño no establece una jerarquía entre usuarios. Para permitir la bidireccionalidad de los mensajes y cumplir con las restricciones de clave ajena en la tabla \texttt{MENSAJE}, se registra la relación en ambos sentidos (\textit{idusuario1, idusuario2} e \textit{idusuario2, idusuario1}). Esto garantiza que cualquier mensaje enviado, independientemente de quién sea el emisor, encuentre una referencia válida en esta tabla.
    \item \textbf{Integridad Referencial (ON DELETE CASCADE):} Esta instancia no es necesaria, pero se incluye por si acaso. Al borrarse de manera lógica los usuarios, se mantiene todo, sin embargo, en la función de listar usuarios, se omiten aquellos borrados, por lo que no se visualizan los mensajes.
\end{itemize}

\subsection{Tabla: MENSAJE}
Almacena el contenido de las interacciones. Cada registro está vinculado a la tabla \texttt{CONVERSA}, lo que asegura que solo se puedan intercambiar mensajes entre usuarios que posean un canal de comunicación previamente validado.

\begin{lstlisting}[language=SQL, caption={Creación de la tabla MENSAJE}]
-- Creación de la tabla mensaje
CREATE TABLE MENSAJE (
    IDMENSAJE           NUMBER GENERATED BY DEFAULT AS IDENTITY,
    IDUSUARIO1          NUMBER NOT NULL,
    IDUSUARIO2          NUMBER NOT NULL,
    MENSAJE             VARCHAR2(1000) NOT NULL,
    FECHAENVIO          DATE DEFAULT SYSDATE NOT NULL,
    BITVISTO            VARCHAR2(1) DEFAULT 'N' NOT NULL,
    
    CONSTRAINT PK_MENSAJE PRIMARY KEY (IDMENSAJE),
    CONSTRAINT FK_MSG FOREIGN KEY (IDUSUARIO1,IDUSUARIO2) REFERENCES CONVERSA(IDUSUARIO1, IDUSUARIO2) ON DELETE CASCADE,
    CONSTRAINT CHECK_N_Y CHECK (BITVISTO IN ('Y','N'))
);
\end{lstlisting}

\textbf{Justificación de diseño:}
\begin{itemize}
    \item \textbf{Referencia Directa a Conversación:} La clave ajena \texttt{FK\_MSG} obliga a que el par emisor-receptor coincida exactamente con una entrada en \texttt{CONVERSA}. Gracias a la naturaleza simétrica del diseño de la tabla superior, se pueden recuperar todos los mensajes de una misma conversación consultando ambas combinaciones de la pareja de usuarios.
    \item \textbf{Restricción CHECK (CHECK\_N\_Y):} Valida que el estado de lectura (\texttt{BITVISTO}) solo admita los valores 'Y' (visto) o 'N' (no visto), garantizando la consistencia de los datos para la lógica de notificaciones.
\end{itemize}

\subsection{Tabla: ARCHIVADO}
Gestiona las preferencias individuales de visualización. Permite que un usuario archive su vista de la conversación sin afectar la disponibilidad de los mensajes para el otro participante.

\begin{lstlisting}[language=SQL, caption={Creación de la tabla ARCHIVADO}]
-- Creación de la tabla archivado
CREATE TABLE ARCHIVADO (
    IDUSUARIO1          NUMBER NOT NULL,
    IDUSUARIO2          NUMBER NOT NULL,
    
    CONSTRAINT PK_ARCHIVADO PRIMARY KEY (IDUSUARIO1, IDUSUARIO2),
    CONSTRAINT FK_ARCH_U1 FOREIGN KEY (IDUSUARIO1) REFERENCES USUARIO(IDUSUARIO) ON DELETE CASCADE,
    CONSTRAINT FK_ARCH_U2 FOREIGN KEY (IDUSUARIO2) REFERENCES USUARIO(IDUSUARIO) ON DELETE CASCADE
);
\end{lstlisting}

\textbf{Justificación de diseño:}
\begin{itemize}
    \item \textbf{Unicidad de Preferencia:} La clave primaria compuesta asegura que un usuario no pueda archivar varias veces al mismo interlocutor. Al ser una tabla independiente de \texttt{CONVERSA}, permite una gestión flexible y desacoplada de la interfaz de usuario.
    \item \textbf{Integridad Referencial (ON DELETE CASCADE):} No es necesario, por el borrado lógico de usuarios, no se elimina.
    
\end{itemize}